\begin{abstract}
Fodor and Pylyshyn argued that thought is systematic, because a mind that can think that John loves Mary can think that Mary loves John, and they traced this ability to representations built from reusable parts. Machine learning usually tests the ability and leaves the representation unscored. Behavioral compositionality is the ability of a system's own procedures to give correct outputs on new combinations of familiar parts. However, behavioral compositionality fails on parts outside training, on additions to a trained system, on composites deeper than training and on the roles of the parts, and a lookup over a finite vocabulary passes many of its tests. Representational compositionality is the ability of a system's states to keep every part of a whole in its role, so that a general procedure can read it back, and prior work ties this property to guarantees of generalization and to the reuse of parts across wholes. We maintain that the reasons for valuing composition, the reuse of parts beyond experience, make representational compositionality the property to seek. We propose recoverability, the return of every part in its role from composites of parts fresh to the reading procedure, as a soft heuristic for designing such systems. We contribute a calibrated metric for recoverability and proofs of its limits. In networks trained on toy domains and in pretrained language, text and image models, we show how behavior fails and that procedures reading recoverable states repair those failures on the same fresh parts. The two properties dissociate in both directions, so a claim of compositionality should say which one it concerns.
\end{abstract}

%=====================================================================
\section{Introduction}
\label{sec:intro}

Compositionality matters because a finite experience must support an unbounded number of new situations. People understand sentences they have never heard by building their meaning from familiar parts \citep{frege1923}, anyone who can think that John loves Mary can think that Mary loves John \citep{fodor1988}, and a person who learns to \emph{dax} can immediately dax twice \citep{lake2018}. Composition earns its value through the abilities that follow from reuse. A compositional system handles new combinations of known parts, takes in a new part and uses it in every role, adds knowledge without damaging old knowledge, nests wholes inside wholes, and keeps each part in its role.

Machine learning tests these abilities mostly through behavior. Benchmarks of compositional generalization score a model's outputs on new combinations of familiar parts \citep{lake2018,keysers2020,kim2020cogs,hupkes2020}, and text-to-image models compose concepts into new scenes \citep{ramesh2021,okawa2023}. This behavior shows some of the benefits of composition and leaves others untested. A lookup over a finite vocabulary passes a test on familiar parts, success can come from fragments seen in training \citep{dziri2023}, and text-to-image models misplace objects \citep{conwell2022} and attach an attribute to the wrong object \citep{rassin2022}. A correct output on one combination says little about new parts, later additions, deeper composites or the roles of the parts.

A second tradition looks inside the system. A classical representation contains its constituents \citep{fodor1988}, and connectionist representations keep constituents that general procedures can read back \citep{smolensky1990,vangelder1990}. Such a representation promises the abilities that behavior leaves untested, because any procedure can reach each part. Identifiability results tie this structure to generalization \citep{lachapelle2023,wiedemer2024,uselis2026}.

We call the first property \emph{behavioral} and the second \emph{representational} compositionality. Both concern the return of parts from a whole. A behavioral test reads the parts through the system's own procedures, and a representational test reads them through a readout that the experimenter declares. We use \emph{recoverability}, the return of each part in its role from composites of parts fresh to the reading procedure, to define the second property and to compare the two on the same parts. Fresh parts are essential, because a readout fitted on a few familiar parts decodes almost any combination of them, blends included. We maintain that the reasons for valuing composition make representational compositionality the property to seek, and we offer recoverability as a soft heuristic for designing such systems. The heuristic is soft because recoverability comes in degrees and depends on a declared class of readouts.

We first show how systems with behavioral compositionality fail, in systems trained from scratch and in pretrained models. We then show that recoverable states support the abilities that behavior lacks, in the same two settings. A recoverable state does not by itself make a system's own procedures read it, and the two properties dissociate in both directions (Table~\ref{tab:twobytwo}). The value of a recoverable state lies in the procedures it makes possible.

\textbf{Contributions.} We define behavioral and representational compositionality so that the two can be measured on the same fresh parts, and we give a calibrated metric $\kappa$ for recoverability. We prove that a lookup decodes almost any relation on a finite vocabulary, that the rank of a linear relation and the width of a state bound the recovery of fresh parts, and that a shared code fixes every role reversal from one composite. We show that behavior fails on fresh parts, on added parts and at depth in networks trained from scratch, and on the roles of the parts in pretrained text, image--text and language models. We show that recoverable states hold the parts these procedures lose, that decoders use the role codes, and that procedures reading through recoverable states repair each failure. In language models with up to 15 billion parameters, we show that token-averaged states keep word order, that the state at the end of a question carries the role, and that this state tracks the model's answers.

%=====================================================================
\section{Related work}
\label{sec:related}

\textbf{Compositionality of meaning.} Frege's principle holds that the meaning of a complex follows from the meanings of its parts and the way they combine \citep{frege1892,frege1923}, and \citet{montague1970} made it a homomorphism from syntax to semantics. As a constraint on functions alone the principle is vacuous, because a recoding of the parts makes any assignment of meanings compositional \citep{zadrozny1994,westerstahl1998}. Conceptual combination strains it from the other side. A guppy is a typical pet fish and neither a typical pet nor a typical fish \citep{osherson1981,fodor1996}, a conjunction can carry attributes that neither concept has \citep{hampton1987}, and \emph{red} contributes differently to \emph{red hair} and \emph{red apple} \citep{kamp1995}.

\textbf{Compositional generalization.} Machine learning studies compositionality through behavior. Benchmarks score a model's outputs on new combinations of familiar parts \citep{lake2018,keysers2020,kim2020cogs}, \citet{hupkes2020} divide the property into systematicity, productivity, substitutivity, localism and overgeneralization, and meta-learning brings networks to the human level on systematic generalization \citep{lake2023}. Models with similar accuracy on familiar data generalize very differently \citep{mccoy2020berts}, success on new combinations can come from matching fragments seen in training \citep{dziri2023}, and agents can generalize without compositional messages \citep{chaabouni2020}. These tests leave the states unscored.

\textbf{Constituent structure.} \citet{fodor1988} argued that the systematicity and productivity of thought require representations whose constituents are tokened in every whole that contains them, and they challenged connectionist models to explain both without implementing a classical architecture. \citet{smolensky1987,smolensky1990} answered with tensor-product representations, whose constituents return by unbinding without appearing explicitly in the state, and \citet{vangelder1990} called such representations functionally compositional, requiring only general procedures that build and decompose them. Recursive auto-associative memories \citep{pollack1990} learned such procedures, holographic reduced representations fixed them \citep{plate1995}, and \citet{chalmers1990} transformed their states without decomposing them. \citet{fodor1990} replied that recoverable constituents need not act in processing, and the dispute over whether systematicity must follow from the architecture continues \citep{aizawa2003,quiltydunn2023,russin2024,woydt2025,goodale2026}. Models of analogy bind fillers to roles explicitly \citep{hummel1997,doumas2022}.

\textbf{Representation and generalization.} Identifiability results guarantee generalization for additive decoders that an encoder inverts \citep{brady2023,lachapelle2023,wiedemer2024} and for linear, orthogonal factors \citep{uselis2026}. Probes decode constituents \citep{bernardi2020,redhardt2025}, yet decodable factors need not yield a decoder that generalizes \citep{montero2021,schott2022}. Tensor-product decompositions \citep{mccoy2019,soulos2020}, binding identifiers \citep{feng2024} and interchange interventions \citep{geiger2021,geiger2024das} can supply part codes, and compression measures score symbolic redescriptions of a representation \citep{andreas2019,elmoznino2024,elmoznino2026}. We add the requirement that a readout transfer to parts outside the vocabulary that fitted it.

%=====================================================================
\section{Behavioral and representational compositionality}
\label{sec:definition}

In this section, we define both properties through the recovery of parts, by the system's own procedures in one case and by a declared readout in the other. We then describe the structure that produces recoverable states and the results that follow. Appendix~\ref{app:proofs} gives the proofs.

@@FIG:fig:ops@@

\subsection{Parts, composites and recovery}
\label{sec:intuition}

A part is a draw from a distribution $\D$ over possible parts, such as glyphs, nouns or photographs. A system assigns each part $x$ a \emph{part code} $e(x) \in \R^d$, and it assigns a composite $r(x_1, \ldots, x_k)$ of $k$ parts in the roles of a relation $r$ a state $h = g_r(x_1, \ldots, x_k) \in \R^m$. Training exposes a system to a finite \emph{vocabulary} $\Voc$ of \emph{familiar} parts. A part is \emph{fresh} to a procedure when no composite containing it helped fit the procedure. A glyph can be fresh both to a readout and to the network that encodes it, while a noun in a pretrained model can be fresh only to the readout. \emph{Recovery} is the return of a part in its role from a composite, and a procedure recovers a fresh part only if it follows the relation beyond its vocabulary.

Suppose a system adds two part vectors $x$ and $y$. Every pair on the line $x + y = z$ yields the same composite (Figure~\ref{fig:ops}a), so no procedure can return the parts. A tensor-product binding $x \otimes \rho_1 + y \otimes \rho_2$ with independent role vectors keeps one decomposition for each composite (Figure~\ref{fig:ops}b). A small vocabulary hides the difference, because few vocabulary pairs share a line and a lookup finds the summed parts. A fresh pair exposes it, because its line holds infinitely many other decompositions.

\subsection{Behavioral compositionality: recovery by the system's own procedures}
\label{sec:notions}

Behavior scores the system's own outputs. A decoder renders parts, a decomposer returns them, and a language model answers questions about them. Each is a recovery of parts through a procedure that the system itself learned.

@@DEF:behcomp@@

This definition asks for fresh parts because familiar parts decode anyway. A non-invertible relation can be injective on a finite vocabulary, and a lookup table then inverts it (Lemma~\ref{lem:lookup}), as clean-up does for vector-symbolic codes \citep{plate2003,frady2018}. Any procedure learned from a vocabulary shares this limit. Composites fix an affine procedure only on the codes they contain, so the procedure covers fresh composites only when the codes of the training parts in each role have full rank $d$ (Proposition~\ref{prop:span}). This result is a procedure-side form of the argument of \citet{marcus1998} that networks fail to extend a relation outside their training space. Behavior on familiar parts therefore reveals little about composition, and a system can fail on fresh parts for two separate reasons, a state that lost them or a procedure that cannot read them.

\subsection{Representational compositionality: recovery by a declared readout}
\label{sec:repcomp}

A representational test separates those two reasons by reading the state directly. A \emph{readout class} $\U$, such as affine maps, contains the procedures allowed to recover a part's code from a state. The class must be fixed in advance, because when the coordinates of an additive decomposition are unrestricted, every invertible relation is additive (Lemma~\ref{lem:psi}).

@@DEF:recoverable@@

@@DEF:repcomp@@

@@DEF:informative@@

Recoverability asks a fixed procedure of a declared kind to return each part in its role across the whole part distribution, and $1 - \varepsilon$ is the population $R^2$ of the best readout in $\U$. Fitting on a large fitting pool and scoring on fresh parts removes the biases of a small vocabulary (Lemma~\ref{lem:lookup} and Proposition~\ref{prop:span}). A relation fails recoverability when it discards information about its parts. For a symmetric relation of identically distributed parts, no readout can tell the roles apart, so in population $\varepsilon_{\mathrm{blend}} = 1/2$, $\varepsilon_{\mathrm{ref}} = 0$ and $\kappa = 1 - 2\varepsilon$. The references correct the estimates for finite samples and readout fitting. Unrestricted readouts let any injective system pass, as unrestricted probes \citep{hewitt2019} and causal alignment maps \citep{sutter2025} do. The informativeness clause closes a second loophole. The blend $(Px, Py)$, with $P$ a projection of rank $r < d$, discards $d - r$ dimensions of each part, yet with the code $e(x) = Px$ it is exactly recoverable. We report informativeness wherever a description of the parts exists apart from the system, which holds for glyph latents and fails for nouns and photographs. The definition depends on $\M$ because only a search can find a decomposition, and a negative result reports that the declared methods found none.

The metric $\kappa$ makes recoverability a soft heuristic. It places a state on a scale from a blend to a role--filler reference, so a designer can compare architectures, pooling rules and checkpoints before any procedure reads the state. Recoverability also puts the two properties on one scale. A declared readout measures the representation, and the system's own procedure measures behavior, scored by the error of Eq.~\ref{eq:recover} when it outputs parts and by accuracy when it chooses among them. Both act on composites of the same fresh parts, so the gap between them dissociates the two properties (Table~\ref{tab:twobytwo}). A system that meets Definitions~\ref{def:behcomp} and~\ref{def:repcomp} is a \emph{compositional system}.

@@TAB:tab:twobytwo@@

\subsection{Shared codes produce recoverable states}
\label{sec:mechanics}

Recoverability describes an end effect, and a state can produce it in several ways. The simplest gives each part one code that enters every whole, changed only by its role.

@@DEF:factor@@

A role-factorized state with full-rank role maps returns every part's code by an affine map, with an error that grows with $\delta$, and one composite then fixes every role reversal (Propositions~\ref{prop:systematic} and~\ref{prop:approx}). Role-factorization is the classical demand that a constituent contribute the same thing to every complex containing it \citep{fodor1988}. The maps $B_i$ must be injective, because otherwise $e'$ could hold a separate code for each role and two unrelated lookup tables would pass. A relation fails role-factorization when a part's contribution depends on its co-parts, as \emph{red} does in \emph{red hair} and \emph{red apple} \citep{kamp1995}. A state that gives each role its own code can be recoverable with a large $\delta$, so $\delta$ diagnoses the mechanism behind recoverability and adds no condition to Definition~\ref{def:repcomp}.

\subsection{Three consequences of recoverability}
\label{sec:consequences}

Appendix~\ref{app:proofs} proves two lemmas and seven propositions, and three results beyond those above carry the argument.

\textbf{Blends lose a computable share of their parts.} A linear relation $g(x_1, \ldots, x_k) = A x$ with $A = [A_1 \cdots A_k]$ gives each composite an affine set of decompositions of dimension $kd - \mathrm{rank}\,A$ (Proposition~\ref{prop:fiber}). For Gaussian parts, no readout of any kind recovers role $i$ with $R^2$ above $\mathrm{tr}(A_i^{\top}(AA^{\top})^{+}A_i)/d$, whose mean over roles is $\mathrm{rank}\,A/(kd)$ (Proposition~\ref{prop:ceiling}). The ceiling is 0.8 for role~1 of $x + 0.5y$, 0.5 for $x + y$, and $m/(kd)$ for a generic projection into $m < kd$ dimensions. For any relation and part distribution, an affine readout of an $m$-dimensional state explains at most the share of the parts' variance in the top $m$ principal components of the stacked part codes (Proposition~\ref{prop:ey}).

\textbf{A shared code fixes reversals from one composite.} If a relation is role-factorized with $\delta = 0$ and $[B_1 \cdots B_k]$ has full column rank, an affine function of a single composite gives every role reversal and every assignment of a known part to a new role (Proposition~\ref{prop:systematic}). This is the unbinding argument for tensor-product representations \citep{smolensky1990}. \citet{fodor1990} objected that recoverable constituents need not act in processing. An interchange intervention that replaces one role's component of a state tests this objection, because it should change the system's behavior for that role alone.

\textbf{Productivity needs invertibility on the relation's own outputs.} A nested linear blend of depth $n$ into a state of width $d$ loses at least $(2^n - 1)d$ dimensions of its leaves (Proposition~\ref{prop:depth}). A relation that inverts on composites of parts can still fail on composites of composites, because they come from another distribution.

%=====================================================================
\section{Experimental design}
\label{sec:experiments}

In this section, we present the protocol that every experiment shares and two designs. Systems trained from scratch let a part be fresh to the system itself, so their own procedures can be scored on parts outside their training. Pretrained models saw their parts in pretraining, so a part is fresh only to the readout, and random weights and training checkpoints supply the controls.

@@FIG:fig:protocol@@

\paragraph{Protocol.} We split the parts before fitting anything. Readouts and role maps fit composites of a fitting pool, validation parts choose ridge penalties, layers and part codes, and fresh parts give every reported number. A ridge readout of the system's state returns each role's part code, and the identical pipeline, applied to the concatenation and to the sum of the same part codes, calibrates its error into $\kappa$ (Eq.~\ref{eq:kappa}). The system's own procedure, whenever it has one, is scored on the same fresh parts. Each system also runs with random weights, which keep its architecture, tokenizer and input format and remove learning, so the gap between trained and random weights measures the contribution of training. Every number is a mean over seeds, and Appendix~\ref{app:protocol} gives further details.

\subsection{Systems trained from scratch}
\label{sec:design_scratch}

\paragraph{Analytic operators.} Fourteen operators calibrate the measure. Parts are Gaussian vectors in $\R^{16}$, and the operators combine them in pairs. They include role--filler binding (concatenation and a tensor product), fixed-width superposition (random projections into $0.5$ to $1.5$ times $kd = 32$ dimensions, HRR and MAP binding, and $x + 0.5y$), symmetric blends ($x + y$ and $x \odot y$), and a random $\tanh$ layer. We fit ridge, least-squares and MLP readouts of the role-1 part on composites of a vocabulary of $V$ parts, with 5 seeds for each $V$.

\paragraph{Glyph networks.} A fixed random generator renders a latent $z \in \R^{8}$ as an $8 \times 8$ glyph, and a composite places two glyphs on an $8 \times 16$ canvas under one of three relations. \textsc{beside} uses disjoint regions, \textsc{over} draws both glyphs in one region with the second mirrored, and \textsc{occlude} draws the second at half scale over the first, a non-additive relation. An encoder writes a state $h \in \R^m$, and a decoder renders the same parts in the same roles under a target relation. The state is therefore the system's representation, and the decoder is its own procedure. We train an MLP, a CNN, a transformer and a role--filler decoder that decodes two halves of $h$ with one shared glyph decoder, at widths $m$ from 8 to 64 and vocabularies $V$ from 8 to 64 glyphs, with 9 seeds. For the tests of role codes, we declare the top 16 principal components of glyph pixels as the part code and retrain each network with 5 seeds. We compare the decoder's rendering of fresh glyphs with a \emph{population readout} fitted on the fitting pool and with a \emph{vocabulary readout} fitted on the network's own training glyphs.

\paragraph{Recursive composers.} Balanced binary trees of depth $n$ have Gaussian leaves in $\R^8$. A recursive auto-associative memory \citep{pollack1990} applies one learned composer at every node and one learned decomposer recursively, and we train both on depths 1 and 2. The decomposer is the system's own procedure, and a population readout maps the same root state to all $2^n$ leaves. Holographic reduced representations (HRR) of the same width $w$ use one pair of unitary role vectors per level, with 5 seeds for each $w$ from 16 to 256.

\subsection{Pretrained models}
\label{sec:design_pretrained}

\paragraph{Parts and states.} Sentences place two nouns in the roles of five asymmetric relations, such as ``a $x$ on a $y$'', and of the control ``a $x$ and a $y$''. We split 305 concrete nouns into 165 for fitting readouts, 40 for validation and 100 fresh to the readout. Photographs come from Tiny-ImageNet \citep{le2015tiny}, whose 200 classes we split into 120 for fitting, 30 for validation and 50 fresh, and two photographs share a $224 \times 224$ canvas \textsc{beside} or \textsc{above} each other, or \textsc{overlap}ping with the second covering a corner of the first. A part code is the model's state for one part alone (``a $x$'', ``there is a $x$ here'', or the photograph), reduced to 16 principal components and whitened, because a few directions dominate transformer states \citep{ethayarajh2019,timkey2021}. The state of a composite is the \emph{token-averaged state} over its tokens or patches. Image states also concatenate the mean patch state of each quadrant, the \emph{quadrant state}, and language models are also read at single tokens. Each split is redrawn for each of 3 seeds.

\paragraph{Models.} We test eleven text encoders, three from CLIP \citep{radford2021clip}, including those that condition Stable Diffusion 1.5 and 2 \citep{rombach2022}, three from T5 \citep{raffel2020t5}, two from BERT \citep{devlin2019bert} and three from GPT-2 \citep{radford2019gpt2}. The language models are Qwen3-8B \citep{yang2025qwen3}, OLMo-2 7B and 13B \citep{olmo2024olmo2}, phi-4 \citep{abdin2024phi4}, and seven Pythia models from 160 million to 12 billion parameters trained on the same data in the same order \citep{biderman2023pythia}. We follow ten checkpoints of Pythia-1.4b from initialization, and nine checkpoints each of Pythia-410m and of five PolyPythia runs that retrain it from other random seeds \citep{vanderwal2025polypythias}. The image encoders are CLIP ViT-B/32 and ViT-H/14 \citep{radford2021clip,cherti2023}, DINOv2 small, large and giant \citep{oquab2024dinov2}, and SigLIP so400m \citep{zhai2023siglip}. Every model also runs with random weights, and Pythia's step 0 serves as its own control. Causal models let the first noun's tokens see only earlier tokens, so we also report $\kappa$ for each role.

\paragraph{Roles by meaning.} Each relation has a \emph{converse} paraphrase that reverses the order of mention of the \emph{canonical} sentence and keeps the meaning, such as ``a tulip on a barrel'' and ``a barrel under a tulip'' (Figure~\ref{fig:protocol}c). The tulip is the \emph{figure} in both, and the barrel is the \emph{ground}. A readout of the figure, fitted on canonical sentences and scored on converse ones, has a \emph{transfer} of 1 if the state codes semantic roles and 0 if it codes the order of mention.

\paragraph{Behavior.} The language models answer questions about the same sentences that name neither noun, such as ``Which object is on top?'', after four demonstrations with fitting nouns. Half of the questions ask for the figure and half for the ground, in canonical and converse phrasings, so a model that answers by word order scores 0.5, and we score whether the model prefers the correct noun on the 100 fresh nouns. The image--text encoders match each composite of fresh photographs against a caption that names its objects in their places and against the caption with the places swapped. The text-only encoders match a sentence against its paraphrase in the other phrasing and against the same paraphrase with the nouns exchanged. Image-only encoders produce no output of their own, so their tests measure only the representation.

